\noindent\textbf{Exercise 1.1.15 (Uniqueness of Jordan measure).}\par

Let $d\geq1$ and let $m':\mathcal J(\mathbb R^d)\to\mathbb R_{\geq0}$ satisfy
non-negativity, finite additivity, and translation invariance.
Show that there exists $c\geq0$ such that
\[
  m'(E)=c\,m(E)\qquad \text{for every }E\in\mathcal J(\mathbb R^d).
\]
In particular, if $m'([0,1)^d)=1$, show that $m'=m$.

\begin{proof}
Set $c:=m'([0,1)^d)\geq 0$.
Since every elementary set is Jordan measurable,
the restriction of $m'$ to $\mathcal E(\mathbb R^d)$
satisfies the three axioms. Thus, by Exercise 1.1.3,
\begin{equation}
  m'(A)=cm(A)
  \qquad\text{for every elementary set }A.
  \tag{1}
\end{equation}

We first note that $m'$ is monotone on
$\mathcal J(\mathbb R^d)$. Indeed, if
$E\subseteq F$ are Jordan measurable, then
$F\setminus E$ is Jordan measurable by
Exercise 1.1.6. Since
$F=E\sqcup(F\setminus E)$, finite additivity
and non-negativity give
\[
  m'(F)=m'(E)+m'(F\setminus E)\geq m'(E).
\]

Now let $E\in\mathcal J(\mathbb R^d)$ be arbitrary.
By the definition of Jordan measurability,
\begin{equation}
 m(E)
 =\sup_{\substack{A\subseteq E\\ A\text{ elementary}}}m(A)
 =\inf_{\substack{E\subseteq B\\ B\text{ elementary}}}m(B).
 \tag{2}
\end{equation}

For every elementary set $A\subseteq E$,
monotonicity and (1) give
\[
  cm(A)=m'(A)\leq m'(E).
\]
Thus $m'(E)$ is an upper bound for all such
$cm(A)$. Taking the supremum and using (2),
together with $c\geq0$, we obtain
\begin{equation}
 cm(E)
 =\sup_{\substack{A\subseteq E\\ A\text{ elementary}}}cm(A)
 \leq m'(E).
 \tag{3}
\end{equation}

Similarly, for every elementary set $B\supseteq E$,
monotonicity and (1) yield
\[
  m'(E)\leq m'(B)=cm(B).
\]
Hence $m'(E)$ is a lower bound for all such
$cm(B)$. Taking the infimum and using (2),
we obtain
\begin{equation}
 m'(E)
 \leq\inf_{\substack{E\subseteq B\\ B\text{ elementary}}}cm(B)
 =cm(E).
 \tag{4}
\end{equation}

Combining (3) and (4), we conclude that
\[
  \boxed{m'(E)=cm(E).}
\]
Since $E$ was arbitrary, the identity holds
for every Jordan measurable set.

In particular, if $m'([0,1)^d)=1$, then
$c=1$, and hence $m'=m$.
\end{proof}
